\section{Parties finies, régularisation et lecteurs résiduels}
\label{lectura:manuscrito-04}
La fonctionnelle construite à partir des modes entiers admet des opérations qui conservent des aspects distincts de sa structure. L’évaluer en un point de convergence, extraire le terme régulier qui accompagne un pôle et dériver son prolongement méromorphe sont trois opérations précises. Cette distinction organise les constantes de ce chapitre et leur relation avec celles du chapitre des moments.

L'antécédent arithmétique est la représentation des formes normales obtenues avec les deux feuillets d'APP. Son évaluation identifie une seule fois chaque entier positif ; l'opérateur de numération agit sur la base correspondante par

\[
Qe_n=ne_n,\qquad
\operatorname{Dom}Q=\left\{u:\sum_{n\geq1}n^2|u_n|^2<\infty\right\}.
\]
La lecture spectrale ultérieure est

\[
Z(s)=\operatorname{Tr}(Q^{-s})=\sum_{n\geq1}n^{-s},\qquad \Re s>1.
\]
Le chapitre de la forme complétée démontre sa coïncidence avec la fonctionnelle thêta–Mellin et son prolongement. On réunit ici une procédure d'évaluation avec des restes contrôlés. La mémoire de l'historique demeure dans l'état source ; la trace est une lecture des modes exprimés, non une reconstruction de tout cet historique à partir d'une constante scalaire.

\subsection{Une formule de prolongement avec reste explicite}
Définissons les nombres de Bernoulli par l’identité formelle

\[
\frac{t}{e^t-1}=\sum_{j\geq0}B_j\frac{t^j}{j!}.
\]
Ses coefficients se calculent sans table numérique :

\[
B_0=1,\qquad
B_m=-\frac1{m+1}\sum_{j=0}^{m-1}\binom{m+1}{j}B_j.
\]
En particulier, \(B_1=-1/2\). Écrivons \((s)_r=s(s+1)\cdots(s+r-1)\), avec \((s)_0=1\), et notons \(B_m(x)=\sum_{j=0}^m\binom mjB_jx^{m-j}\) le polynôme correspondant. Sa périodisation est \(\widetilde B_{2p}(x)=B_{2p}(\{x\})\). Pour \(N,p\geq1\), la sommation d’Euler–Maclaurin donne

\[
\begin{aligned}
Z(s)={}&\sum_{n=1}^N n^{-s}+\frac{N^{1-s}}{s-1}-\frac1{2N^s}\\
&+\sum_{k=1}^p\frac{B_{2k}}{(2k)!}(s)_{2k-1}N^{1-s-2k}
+R_{N,p}(s),
\end{aligned}
\]
avec

\[
R_{N,p}(s)=-\frac{(s)_{2p}}{(2p)!}
\int_N^\infty\widetilde B_{2p}(x)x^{-s-2p}\,dx.
\]
\HMTresultado{Proposition}{Domaine et borne du reste}{rz-num-04-1}
L'intégrale du reste est holomorphe dans \(\Re s>1-2p\). Si \(\sigma=\Re s>1-2p\), alors

\[
|R_{N,p}(s)|\leq
\frac{|B_{2p}|\,|(s)_{2p}|}{(2p)!}
\frac{N^{1-\sigma-2p}}{\sigma+2p-1}.
\]
La formule précédente prolonge \(Z\) à ce demi-plan, avec son pôle simple en \(s=1\), et elle est compatible avec le prolongement thêta–Mellin.

\textbf{Démonstration.} La formule de sommation s’obtient en intégrant successivement par parties dans chaque intervalle unitaire et en utilisant \(B_j'(x)=jB_{j-1}(x)\) et les égalités aux extrémités pour les polynômes de degré différent de un. En l’appliquant à \(x^{-s}\) sur \([N,M]\), les dérivées fournissent les facteurs \((s)_r\). Dans \(\Re s>1\), le passage \(M\to\infty\) élimine les extrémités supérieures et produit exactement l’identité et le signe du reste écrits ci-dessus.

Pour contrôler celui-ci, on utilise

\[
|B_{2p}(t)|\leq |B_{2p}|,\qquad 0\leq t\leq1.
\]
Une preuve de cette inégalité est la série de Fourier absolument convergente

\[
B_{2p}(t)=(-1)^{p+1}\frac{2(2p)!}{(2\pi)^{2p}}
\sum_{m\geq1}\frac{\cos(2\pi mt)}{m^{2p}}.
\]
Le maximum de la borne par valeurs absolues est atteint en \(t=0\). Cette identité se déduit des coefficients de Fourier par intégration par parties des polynômes eux-mêmes ; le facteur \(\pi\) exprime seulement la convention de Fourier de cette preuve et n’entre pas numériquement dans l’évaluateur rationnel qui sera employé ensuite. La comparaison avec \(x^{-\sigma-2p}\) démontre la borne du reste. Sur les compacts du demi-plan, elle domine aussi ses dérivées en \(s\), qui ajoutent des puissances de \(\log x\) ; l’intégrale y est donc holomorphe. L’identité dans le demi-plan de convergence et l’unicité du prolongement prouvent la compatibilité annoncée. \hfill\(\square\)

L'ordre \(p\) et l'horizon \(N\) sont des paramètres de calcul et de contrôle de l'erreur. Ils ne sélectionnent pas une constante différente : les expressions correspondant à différents \(N,p\) représentent la même fonctionnelle sur leurs domaines communs. Il n'est pas non plus nécessaire qu'une série asymptotique converge lorsque \(p\) augmente indéfiniment avec \(N\) fixé. Chaque paire finie est accompagnée de sa propre borne démontrée.

\subsection{Euler–Mascheroni comme partie finie de la même fonctionnelle}
La partie finie au pôle est définie par

\[
\gamma_0=\lim_{s\to1}\left(Z(s)-\frac1{s-1}\right).
\]
En posant \(s=1+z\), le terme singulier vérifie

\[
\frac{N^{-z}}z=\frac1z-\log N+O(z).
\]
La proposition précédente permet de prendre la partie finie terme à terme.

\HMTresultado{Théorème}{Intervalle explicite pour la partie finie}{rz-num-04-2}
Soient

\[
G_{N,p}=H_N-\log N-\frac1{2N}
+\sum_{k=1}^p\frac{B_{2k}}{2kN^{2k}},
\qquad H_N=\sum_{n=1}^N\frac1n.
\]
Alors

\[
|\gamma_0-G_{N,p}|\leq\frac{|B_{2p}|}{2pN^{2p}}.
\]
De plus,

\[
\gamma_0=\lim_{N\to\infty}(H_N-\log N).
\]
\textbf{Démonstration.} En \(s=1\), on a \((1)_{2k-1}=(2k-1)!\) ; c’est pourquoi chaque terme de Bernoulli devient \(B_{2k}/(2kN^{2k})\). La borne de la proposition~\ref{rz-num-04-1} évaluée en ce point est \(|B_{2p}|/(2pN^{2p})\). Avec \(p=1\), tous les termes qui séparent \(G_{N,1}\) de \(H_N-\log N\), ainsi que le reste, tendent vers zéro. Cela démontre la limite. Il s’agit de l’identification ultérieure à la constante d’Euler–Mascheroni. \hfill\(\square\)

L'antécédent de \(\gamma_0\) utilisé par la formule de Meissel–Mertens est ainsi incorporé. Le nombre n'est pas reçu comme un chiffre indépendant destiné à corriger la somme sur les irréductibles : c'est une opération spécifiée sur la même fonction de partition des entiers.

Le logarithme utilisé dans le calcul admet également un intervalle rationnel. Pour \(x\geq1\), soit \(z=(x-1)/(x+1)\). La série intégrée d’une progression géométrique donne

\[
\log x=2\sum_{j=0}^{M-1}\frac{z^{2j+1}}{2j+1}+r_M(x),
\quad
0\leq r_M(x)\leq
\frac{2z^{2M+1}}{(2M+1)(1-z^2)}.
\]
L'inégalité résulte du remplacement des dénominateurs de la queue par \(2M+1\) et de la sommation de la progression géométrique. On réduit d'abord \(x=2^a y\), avec \(1\leq y<2\), et on évalue \(a\log2+\log y\). Ainsi, l'argument des séries est au plus égal à \(1/3\). Pour \(0<x<1\), on utilise \(\log x=-\log(1/x)\), en inversant les extrémités de l'intervalle.

\subsection{Apéry et l’évaluation négative régularisée}
Pour un entier \(r\geq2\), la spécialisation de la proposition~\ref{rz-num-04-1} produit un centre rationnel

\[
Z_{N,p}^{[r]}=
\sum_{n=1}^N\frac1{n^r}+\frac1{(r-1)N^{r-1}}-\frac1{2N^r}
+\sum_{k=1}^p\frac{B_{2k}(r)_{2k-1}}{(2k)!N^{r+2k-1}}
\]
et la borne entièrement rationnelle

\[
\left|Z(r)-Z_{N,p}^{[r]}\right|
\leq\frac{|B_{2p}|(r)_{2p-1}}{(2p)!N^{r+2p-1}}.
\]
La simplification utilise \((r)_{2p}/(r+2p-1)=(r)_{2p-1}\). En \(r=3\), on obtient la constante d’Apéry. Le procédé détermine sa valeur avec des bornes ; l’irrationalité de cette constante est une assertion mathématique distincte de cette évaluation.

Du côté négatif, le même prolongement admet une évaluation exacte, sans approximation décimale ni paramètre d’ajustement.

\HMTresultado{Théorème}{Évaluation régularisée en moins un}{rz-num-04-3}
Le prolongement de la fonctionnelle des modes entiers satisfait

\[
Z(-1)=-\frac1{12}.
\]
\textbf{Démonstration.} Dans la formule de la proposition~\ref{rz-num-04-1}, on prend \(p=2\), de sorte que \(s=-1\) appartient à l'intérieur du demi-plan de validité. Le facteur \((s)_4\) s'annule en ce point ; l'intégrale qui l'accompagne y est holomorphe, si bien que le reste vaut zéro. S'annule également \((s)_3\), qui multiplie le terme contenant \(B_4\). Il ne reste que

\[
Z(-1)=\sum_{n=1}^N n-\frac{N^2}{2}-\frac N2-\frac{B_2}{2}.
\]
Comme l’accumulation additive donne \(\sum_{n=1}^N n=N(N+1)/2\) et la récurrence de Bernoulli donne \(B_2=1/6\), le résultat est \(-1/12\), indépendamment de \(N\). \hfill\(\square\)

L’objet évalué est le prolongement de la fonctionnelle, non la limite ordinaire des sommes partielles positives, qui croissent sans borne. La distinction conserve l’opération qui fait apparaître la valeur : le terme de Bernoulli demeure après l’annulation des termes polynomiaux de frontière. C’est pourquoi la régularisation possède ici une structure et une généalogie, en plus d’une valeur.

\subsection{Coefficients de Laurent et conservation des trois classes}
Écrivons

\[
Z(1+z)-\frac1z=\sum_{n\geq0}a_nz^n,
\qquad \gamma_n=(-1)^nn!a_n.
\]
La fonction régulière de gauche est holomorphe autour de zéro. Pour un cercle contenu dans son domaine d’holomorphie, Cauchy détermine

\[
a_n=\frac1{2\pi i}\int_{|z|=\rho}
\frac{Z(1+z)-1/z}{z^{n+1}}\,dz.
\]
La formule de prolongement avec reste permet de contrôler cette extraction : une borne uniforme \(\varepsilon\) de l’erreur sur le cercle produit une borne \(\varepsilon\rho^{-n}\) pour le coefficient. Cela justifie un procédé à un ordre arbitraire ; cela ne transforme pas la coïncidence de deux quadratures finies en une borne d’erreur.

La lecture résiduelle maintient les trois classes de l'opérateur entier :

\[
Q_re_n=\mathbf1_{n\equiv r\ (3)}e_n,\qquad
Z_r(s)=\operatorname{Tr}(Q^{-s}Q_r),\quad r=0,1,2.
\]
Ces projections ont un domaine commun dans le demi-plan de convergence. Leur somme est l'identité ; la différence des classes 1 et 2 conserve l'orientation du caractère modulo trois. On a

\[
Z=Z_0+Z_1+Z_2,\qquad Z_0(s)=3^{-s}Z(s),\qquad
L_{-3}(s)=Z_1(s)-Z_2(s).
\]
La dernière série a des coefficients périodiques \((1,-1,0)\) et une moyenne nulle. Le regroupement par périodes permet son prolongement près de \(s=1\). Chaque \(Z_r\) a pour résidu \(1/3\) en ce point.

Soient

\[
Z_r(1+z)-\frac1{3z}=\sum_{n\geq0}c_{r,n}z^n,
\qquad \lambda=\log3.
\]
La comparaison des coefficients donne

\[
a_n=c_{0,n}+c_{1,n}+c_{2,n},\qquad
\frac{L_{-3}^{(n)}(1)}{n!}=c_{1,n}-c_{2,n},
\]
\[
c_{0,n}=\frac13\left[
\frac{(-\lambda)^{n+1}}{(n+1)!}
+\sum_{k=0}^na_k\frac{(-\lambda)^{n-k}}{(n-k)!}\right].
\]
La troisième identité provient de la multiplication de la série complète, y compris son pôle, par \(3^{-1}e^{-\lambda z}\). La somme et la différence récupèrent les deux autres coefficients de manière unique. Les trois lectures conservent, respectivement, l'agrégation, l'orientation et le changement d'échelle.

À l’ordre zéro, on peut achever l’évaluation sans recevoir de décimale de \(L_{-3}(1)\). En effet,

\[
L_{-3}(1)=\sum_{m\geq0}\left(\frac1{3m+1}-\frac1{3m+2}\right)
=\int_0^1\frac{1-x}{1-x^3}\,dx
=\int_0^1\frac{dx}{1+x+x^2}.
\]
On complète le carré au dénominateur et on évalue la primitive :

\[
L_{-3}(1)=\frac2{\sqrt3}
\left[\arctan\frac{2x+1}{\sqrt3}\right]_0^1
=\frac{\pi}{3\sqrt3}.
\]
La coordonnée \(\pi\) est la représentation régionale préalable de HMT ; cette identité est une évaluation analytique ultérieure d'un lecteur de classes, et non une procédure permettant de sélectionner rétroactivement cette coordonnée. Par conséquent, avec \(C_r=c_{r,0}\),

\[
C_0=\frac{\gamma_0-\log3}{3},\qquad
C_{1,2}=\frac{\gamma_0}{3}+\frac{\log3}{6}
\mathbin{\pm}\frac{\pi}{6\sqrt3}.
\]
On récupère les identités

\[
C_0+C_1+C_2=\gamma_0,\qquad
C_1-C_2=L_{-3}(1),\qquad C_1+C_2-2C_0=\log3.
\]
\clearpage
\subsection{Représentation exacte et portée du certificat}
Le programme \texttt{partes\_\allowbreak{}finitas\_\allowbreak{}exactas.\allowbreak{}py} conserve l'ordre suivant : succession de suites au moyen du feuillet additif ; évaluation des modes entiers ; construction rationnelle de Bernoulli ; application de la formule de la fonctionnelle ; incorporation du reste ; représentation décimale dirigée. La multiplication native demeure disponible dans le même ensemble documentaire et sa compatibilité est démontrée dans le chapitre du produit. Le calcul analytique sur les modes exprimés ne modifie pas leurs règles.

Avec \(N=81\), \(p=8\) et 120 termes pour chaque série du logarithme, on obtient les intervalles

\[
\begin{gathered}
0.577215664901532860606512090082273204041535\\
<\gamma_0<\\
0.577215664901532860606512090082531387279784,
\end{gathered}
\]
\[
\begin{gathered}
1.202056903159594285399738161511447311355829\\
<Z(3)<\\
1.202056903159594285399738161511452663119342.
\end{gathered}
\]
Les extrémités décimales sont des arrondis vers l’extérieur de fractions exactes. Le programme conserve également la largeur rationnelle de chaque intervalle et retrouve \(-1/12\) exactement. Aucun chiffre de comparaison ne figure dans le producteur. Le vérificateur séparé ouvre les valeurs de référence après avoir obtenu les sorties ; cette comparaison vérifie l’exécution et ne remplace pas les bornes démontrées.

Ce chapitre clôt dans son domaine les parties finies et les évaluations qu'il vient de définir. Il n'utilise aucune hypothèse sur la localisation des zéros ni sur la positivité de la forme de Weil. Le prolongement vers le problème spectral conserve les obligations propres à ses opérateurs : une évaluation régularisée ne fournit pas à elle seule l'identité de frontière d'une autre fonctionnelle.

Provenance : développement récupéré des fonctionnelles triadiques et du réseau coordonné de lecteurs. La réunion de la preuve d’Euler–Maclaurin, de ses bornes rationnelles et de son programme est une formalisation et un certificat supplémentaire de cette édition ; aucune nouvelle priorité historique n’est attribuée à ces identités analytiques.
